\section{朴素实现为何失败}

在进入 harness 设计之前，作者首先剖析了``单 agent 直接执行''模式的两个系统性失败模式。

\subsection{上下文窗口管理问题}

模型在执行长任务时，随着 context window 逐渐填满，输出的连贯性会显著下降。更微妙的是，Claude 会表现出一种被称为 \textbf{context anxiety} 的行为：当模型``感知到''自己接近上下文长度限制时，会提前草草收尾。

\begin{knowledgebox}{Context Reset vs. Compaction}
作者对比了两种应对上下文膨胀的策略：

\textbf{Compaction（压缩）}：对对话历史进行摘要，缩短后继续。保留了连续性，但无法解决 context anxiety——因为 agent 仍然``记得''自己已经工作了很久。

\textbf{Context Reset（重置）}：彻底清空上下文，启动全新 agent，配合结构化 handoff artifact 传递状态。提供了一个``干净的开始''，但增加了编排复杂度和 token 开销。

在早期测试中，Claude Sonnet 4.5 的 context anxiety 严重到 compaction 不足以应对，因此 context reset 成为 harness 设计的必要组件。
\end{knowledgebox}

\subsection{自我评估偏差}

\begin{warningbox}{Agent 的自我评估陷阱}
\textit{``When asked to evaluate work they've produced, agents tend to respond by confidently praising the work—even when, to a human observer, the quality is obviously mediocre.''}

当要求 agent 评估自己产出的工作时，它们倾向于自信地给出正面评价——即使在人类观察者看来质量明显平庸。这个问题在主观任务（如设计）上尤为突出，因为不存在像软件测试那样的二值化验证手段。
\end{warningbox}

关键洞察在于：\textbf{将生成和评估分离}后，调优一个持怀疑态度的 evaluator 远比让 generator 对自身工作保持批判性更容易。一旦外部反馈存在，generator 就有了具体的迭代目标。

\subsection{本章小结}

朴素的单 agent 实现面临两个结构性困境：（1）上下文窗口膨胀导致的连贯性退化和 context anxiety；（2）自我评估偏差导致的质量停滞。这两个问题共同构成了引入 multi-agent harness 的动机。

